% Part: proof-theory
% Chapter: normalization
% Section: reductions

\documentclass[../../../include/open-logic-section]{subfiles}

\begin{document}

\olfileid{pt}{nor}{red}

\olsection{Konversi Reduksi}

Yèn cut ing !!a{proof} \Log{N1i} dawané~$1$, cut mau segmen
kang mung ngemot siji kedadéan !!{formula}; !!{formula} iku
sekaligus konklusi inferènsi !!a{introduction} lan premis utama
inferènsi !!a{elimination}. Ing bukti normalisasi, cut kaya iki
``direduksi'' kanthi ngganti sub-!!{proof} kang pungkasané pasangan
!!{introduction}/!!{elimination} mau nganggo !!{proof} anyar
kang ora ngemot pasangan inferènsi iku. Panggantian iki bisa
ngasilaké cut anyar. Nanging, cut anyar mau nduwèni pangkat
luwih cilik tinimbang cut kang direduksi. Mula, !!{proof} anyar
nduwèni dawa cut luwih cilik (amarga cut maksimal dawa~$1$
wis dibusak), utawa pangkat cut !!{proof} suda (yèn cut maksimal
kang dimaksud iku siji-sijiné cut maksimal).

Tuladhané, yèn inferènsi kang melu yaiku \Intro{\lif} lan
\Elim{\lif}, $!A \ident !B \lif !C$, lan ana
sub-!!{proof}~$\delta_1$ kang pungkasané mangkéné:
\[
\AxiomC{$\Discharge{!B}{x}$}
\RightLabel{$\delta_2$}
\DeduceC{$!C$}
\DischargeRule{\Intro{\lif}}{x}
\UnaryInfC{$!B \lif !C$}
\AxiomC{}
\RightLabel{$\delta_3$}
\DeduceC{$!B$}
\RightLabel{\Elim{\lif}}
\BinaryInfC{$!C$}
\DisplayProof
\]
Kanggo mbusak cut iki, subbukti~$\delta_1$ ing~$\delta$
diganti mangkéné:
\[
\AxiomC{}
\RightLabel{$\delta_3$}
\DeduceC{$!B$}
\RightLabel{$\Subst{\delta_2}{\delta_3}{!B^x}$}
\DeduceC{$!C$}
\DisplayProof
\]
Kanthi $\Subst{\delta_2}{\delta_3}{!B^x}$ kita maksudaké !!{proof}
asil saka~$\delta_2$ nalika saben asumsi~$!B$ kang
mawa label~$x$ diganti kabèh !!{proof}~$\delta_3$. !!{proof} iki
ora nduwèni manèh asumsi mbukak awujud~$!B$ mawa label~$x$.
Ora ana asumsi liya mawa label~$x$, amarga kaya lumrahé kita
nganggep ora ana rong !!{formula}s kang dadi asumsi mawa label padha.
Mula, !!{proof} anyar nduwèni asumsi mbukak kang padha karo~$\delta_1$;
saben inferènsi ing~$\delta_2$ kang mbusak asumsi isih mbusak
asumsi kang padha ing~$\Subst{\delta_2}{\delta_3}{!B^x}$;
lan saben inferènsi ing~$\delta_3$ kang mbusak asumsi bisa
ngasilaké sawetara salinan kang mbusak asumsi cocog
ing~$\Subst{\delta_2}{\delta_3}{!B^x}$.

Saben cut kang kabèh ana ing $\delta_2$ utawa $\delta_3$,
utawa kabèh ana ing sanjabané~$\delta_1$ ing~$\delta$, ora owah:
cut mau isih nglibataké !!{formula}s kang padha lan ngliwati
inferènsi kang padha, mula pangkat lan dawané padha karo cut
kang cocog ing~$\delta$. Kita wis mbusak cut maksimal dawa~$1$,
mula dawa cut utawa pangkat cut suda (yèn dawa cut~$\delta$
mauné~$1$ lan iki siji-sijiné cut maksimal).

Nanging, ana rong prakara kang bisa kelakon:
\begin{enumerate}
  \item Kanthi ngganti kabèh asumsi $!B^x$ nganggo $\delta_3$,
  kita uga ndadèkaké kabèh cut ing~$\delta_3$ dadi luwih akèh.
  Yèn sawatara iku cut maksimal, dawa cut !!{proof} anyar bisa
  luwih gedhé tinimbang dawa cut !!{proof} lawas. Iki kudu
  dicegah. Mula konversi reduksi mung kita trapaké marang cut
  maksimal kang \emph{paling tengen}, yaiku cut kang ora nduwèni
  cut maksimal liya ing cabang~$\delta$ ing sisih tengené.
  Yèn ana cut ing~$\delta_3$, cut iku ana ing sisih tengené
  cut kang lagi direduksi, mula cut kang dipilih dudu kang paling
  tengen. Kanthi tansah milih cut paling tengen, kita njamin
  yèn dawa cut !!{proof} suda, utawa malah pangkat cuté suda.
  \item Yèn asumsi $!B^x$ ing $\delta_2$ dadi premis utama
  inferènsi !!a{elimination}, lan inferènsi pungkasan~$\delta_3$
  iku \Elim{\lor}, \Elim{\lexists}, utawa inferènsi
  !!a{introduction}, nèmpèlaké~$\delta_3$ menyang $\delta_2$
  ing asumsi mau ngasilaké cut anyar. Kang mauné mung asumsi
  ing~$\delta_2$ saiki dadi cut dawa~$1$ (konklusi inferènsi
  !!a{introduction} lan premis utama inferènsi
  !!a{elimination}), utawa dadi kedadéan !!{formula} pungkasan
  saka segmen cut (premis utama inferènsi !!a{elimination}
  lan konklusi saka~\Elim{\lor} utawa \Elim{\lexists}).
  Yèn $!B^x$ mauné premis tambahan \Elim{\lor} utawa
  \Elim{\lexists}, formula iku wis dadi kedadéan !!{formula}
  kapisan saka sawijining segmen, lan bisa uga segmen cut.
  Ing !!{proof} anyar, segmen iku bisa luwih dawa. Nanging,
  !!{formula} kang mbentuk cut anyar utawa segmen cut kang wis
  ana iku~$!B$, lan $\depth{!B} < \depth{!B \lif !C}$.
  Mula, sanajan cut bisa nambah utawa dawané mundhak,
  ora ana cut anyar mau kang maksimal. Dadi pangkat cut suda,
  utawa, yèn isih ana cut kanthi pangkat padha, dawa cut suda.
\end{enumerate}

Proses nyuda cut dawa~$1$ diarani \emph{konversi reduksi}
(utawa \emph{konversi liku-liku}). Pola kanggo sawetara pasangan
aturan kacetha ing~\olref{tab:reduction-conversions}.
(Pasangan liyané dadi latihan.) 

\begin{table}
\begin{multline*}
\bottomAlignProof
\AxiomC{$\Discharge{!B}{x}$}
\RightLabel{$\delta_2$}
\shortDeduce
\DeduceC{$!C$}
\DischargeRule{\Intro{\lif}}{x}
\UnaryInfC{$!B \lif !C$}
\AxiomC{}
\RightLabel{$\delta_3$}
\shortDeduce
\DeduceC{$!B$}
\RightLabel{\Elim{\lif}}
\BinaryInfC{$!C$}
\DisplayProof
\quad \longrightarrow\quad
\shoveright{\bottomAlignProof
\AxiomC{}
\RightLabel{$\delta_3$}
\shortDeduce
\DeduceC{$!B$}
\RightLabel{$\Subst{\delta_2}{\delta_3}{!B^x}$}
\shortDeduce
\DeduceC{$!C$}
\DisplayProof}
\\[4ex]
\shoveleft{\bottomAlignProof
\AxiomC{}
\RightLabel{$\delta_2$}
\shortDeduce
\DeduceC{$!B$}
\AxiomC{}
\RightLabel{$\delta_3$}
\shortDeduce
\DeduceC{$!C$}
\RightLabel{\Intro{\land}}
\BinaryInfC{$!B \land !C$}
\RightLabel{\Elim{\land}[1]}
\UnaryInfC{$!B$}
\DisplayProof
\quad \longrightarrow\quad
\bottomAlignProof
\AxiomC{}
\RightLabel{$\delta_2$}
\shortDeduce
\DeduceC{$!B$}
\DisplayProof}\qquad
\shoveright{\bottomAlignProof
\AxiomC{}
\RightLabel{$\delta_2$}
\shortDeduce
\DeduceC{$!B$}
\AxiomC{}
\RightLabel{$\delta_3$}
\shortDeduce
\DeduceC{$!C$}
\RightLabel{\Intro{\land}}
\BinaryInfC{$!B \land !C$}
\RightLabel{\Elim{\land}[2]}
\UnaryInfC{$!C$}
\DisplayProof
\quad \longrightarrow\quad
\bottomAlignProof
\AxiomC{}
\RightLabel{$\delta_3$}
\shortDeduce
\DeduceC{$!C$}
\DisplayProof}
\\[4ex]
\bottomAlignProof
\AxiomC{}
\RightLabel{$\delta_2$}
\shortDeduce
\DeduceC{$!B$}
\RightLabel{\Intro{\lor}[1]}
\UnaryInfC{$!B \lor !C$}
\AxiomC{$\Discharge{!B}{x}$}
\RightLabel{$\delta_3$}
\shortDeduce
\DeduceC{$!D$}
\AxiomC{$\Discharge{!C}{y}$}
\RightLabel{$\delta_4$}
\shortDeduce
\DeduceC{$!D$}
\DischargeRule{\Elim{\lor}}{x\,y}
\TrinaryInfC{$!D$}
\DisplayProof
\quad \longrightarrow\quad
\shoveright{\bottomAlignProof
\AxiomC{}
\RightLabel{$\delta_2$}
\shortDeduce
\DeduceC{$!B$}
\RightLabel{$\Subst{\delta_3}{\delta_2}{!B^x}$}
\shortDeduce
\DeduceC{$!D$}
\DisplayProof}
\\[2ex]
\shoveleft{\bottomAlignProof
\AxiomC{}
\RightLabel{$\delta_2$}
\shortDeduce
\DeduceC{$!B(c)$}
\RightLabel{\Intro{\lforall}}
\UnaryInfC{$\lforall[x][!B(x)]$}
\RightLabel{\Elim{\lforall}}
\UnaryInfC{$!B(t)$}
\DisplayProof
\quad \longrightarrow\quad
\bottomAlignProof
\AxiomC{}
\RightLabel{$\Subst{\delta_2}{t}{c}$}
\shortDeduce
\DeduceC{$!B(t)$}
\DisplayProof}
\\[4ex]
\shoveright{\bottomAlignProof
\AxiomC{}
\RightLabel{$\delta_2$}
\shortDeduce
\DeduceC{$!B(t)$}
\RightLabel{\Intro{\lexists}}
\UnaryInfC{$\lexists[x][!B(x)]$}
\AxiomC{$\Discharge{!B(c)}{x}$}
\RightLabel{$\delta_3$}
\shortDeduce
\DeduceC{$!D$}
\DischargeRule{\Elim{\lexists}}{x}
\BinaryInfC{$!D$}
\DisplayProof
\quad \longrightarrow\quad
\bottomAlignProof
\AxiomC{}
\RightLabel{$\delta_2$}
\shortDeduce
\DeduceC{$!B(t)$}
\RightLabel{$\Subst{\Subst{\delta_3}{t}{c}}{\delta_2}{!B(t)^x}$}
\shortDeduce
\DeduceC{$!D$}
\DisplayProof}
\end{multline*}
\caption{Sawetara konversi reduksi kanggo \Log{N1i}.}
\ollabel{tab:reduction-conversions}
\end{table}

\begin{prob}
Wènèhana konversi reduksi kanggo \Intro{\lnot} kang disusul
\Elim{\lnot}.
\end{prob}

Isih ana siji kasus kang kudu dipriksa. Definisi cut uga
nyakup cut kang dawané~$1$ lan $!A$ dadi konklusi
saka~\FalseInt{} uga premis utama inferènsi
!!a{elimination}. Kasus iki ditangani kanthi cara padha
kaya nalika $!A$ dadi konklusi aturan !!a{introduction}.
Tuladhané, gantenana
\[
\bottomAlignProof
\AxiomC{}
\RightLabel{$\delta_2$}
\DeduceC{$\lfalse$}
\UnaryInfC{$!B \lif !C$}
\AxiomC{}
\RightLabel{$\delta_3$}
\DeduceC{$!B$}
\RightLabel{\Elim{\lif}}
\BinaryInfC{$!C$}
\DisplayProof
\quad\text{dadi}\quad
\bottomAlignProof
\AxiomC{}
\RightLabel{$\delta_2$}
\DeduceC{$\lfalse$}
\UnaryInfC{$!C$}
\DisplayProof
\]
Kita kudu luwih ngati-ati yèn $!A \ident (!B \lor
!C)$ utawa $!A \ident \lexists[x][!B(x)]$.
Tuladhané, yèn kita ngganti
\[
\bottomAlignProof
\AxiomC{}
\RightLabel{$\delta_2$}
\DeduceC{$\lfalse$}
\RightLabel{\FalseInt}
\UnaryInfC{$!B \lor !C$}
\AxiomC{$\Discharge{!B}{x}$}
\RightLabel{$\delta_3$}
\DeduceC{$!D$}
\AxiomC{$\Discharge{!C}{y}$}
\RightLabel{$\delta_4$}
\DeduceC{$!D$}
\DischargeRule{\Elim{\lor}}{x\,y}
\TrinaryInfC{$!D$}
\DisplayProof
\quad\text{dadi}\quad
\bottomAlignProof
\AxiomC{}
\RightLabel{$\delta_2$}
\DeduceC{$\lfalse$}
\RightLabel{\FalseInt}
\UnaryInfC{$!D$}
\DisplayProof
\]
kita bisa waé ngasilaké cut anyar kang !!{formula} cuté~$!D$,
lan ora ana jaminan yèn $\depth{!D} < \depth{!B \lor !C}$!{}
Mula kita kudu nindakaké cara kang padha karo reduksi
\Intro{\lor}/\Elim{\lor}, yaiku ngganti nganggo
\[
\bottomAlignProof
\AxiomC{}
\RightLabel{$\delta_2$}
\DeduceC{$\lfalse$}
\RightLabel{\FalseInt}
\UnaryInfC{$!B$}
\RightLabel{$\delta_3$}
\DeduceC{$!D$}
\DisplayProof
\quad\text{utawa}\quad
\bottomAlignProof
\AxiomC{}
\RightLabel{$\delta_2$}
\DeduceC{$\lfalse$}
\RightLabel{\FalseInt}
\UnaryInfC{$!C$}
\RightLabel{$\delta_4$}
\DeduceC{$!D$}
\DisplayProof
\]

\end{document}
